\documentclass[12pt,a4paper]{article}
\usepackage[spanish]{babel}
\usepackage[utf8]{inputenc}
\usepackage[T1]{fontenc}
\usepackage{geometry}
\usepackage{longtable}
\usepackage{array}
\usepackage{hyperref}
\usepackage{setspace}
\geometry{margin=2.5cm}
\setstretch{1.15}

\begin{document}

\begin{titlepage}
    \centering
    {\Large Instituto Tecnologico de Costa Rica\par}
    \vspace{1.5cm}
    {\LARGE \textbf{Proyecto 3: SiembraTEC}\par}
    \vspace{0.4cm}
    {\large Videojuego de granja en Python y Tkinter\par}
    \vspace{2.0cm}

    \begin{flushleft}
        \textbf{Curso:} Programacion Orientada a Objetos\\
        \textbf{Tipo de proyecto:} Individual\\
        \textbf{Autora:} Joselyn Hidalgo Torres\\
        \textbf{Documento:} Documentacion tecnica y manual de usuario
    \end{flushleft}

    \vfill
    {\large \today\par}
\end{titlepage}

\tableofcontents
\newpage

\section{Descripcion del sistema}
El sistema desarrollado corresponde a una simulacion de granja inspirada en FarmVille, construida con Python y Tkinter bajo un enfoque de Programacion Orientada a Objetos (POO). El objetivo del juego es administrar recursos, producir monedas, comprar y colocar elementos productivos, ampliar terrenos y conservar el progreso entre ejecuciones.

La partida inicia con 100 monedas y un terreno base de 10x10 celdas (segun la aclaracion oficial del proyecto). Cada celda admite un elemento o permanece vacia. El jugador puede adquirir plantaciones, animales, arboles, decorativos y consumibles de cura. La produccion genera ganancias por tiempo y alimenta las estadisticas de la partida.

El sistema integra:
\begin{itemize}
    \item Persistencia completa mediante archivo JSON.
    \item Inventario para almacenar elementos comprados o retirados del terreno.
    \item Bitacora para registrar eventos de juego.
    \item Temporizadores con hilos por parcela productiva.
    \item Auto-recoleccion opcional con penalizacion por no recolectar en modo manual.
    \item Sistema aleatorio de plagas y enfermedades con uso de pesticida/medicina.
\end{itemize}

\section{Funcionalidades implementadas}
\subsection{Menu principal}
\begin{itemize}
    \item Iniciar partida nueva.
    \item Confirmacion fuerte para borrar progreso previo al iniciar nueva partida.
    \item Continuar partida guardada.
    \item Ver estadisticas globales.
    \item Salida segura de la aplicacion.
\end{itemize}

\subsection{Terrenos y representacion matricial}
\begin{itemize}
    \item Cada terreno utiliza una matriz de 10x10 (lista de listas).
    \item Validacion de posiciones para evitar filas/columnas invalidas.
    \item Celda libre requerida para colocar nuevos elementos.
    \item Compra de terrenos adicionales por 20.000 monedas.
    \item Navegacion por terreno anterior, siguiente o salto directo por numero.
\end{itemize}

\subsection{Tienda}
\begin{itemize}
    \item Catalogo completo de productos solicitados por el enunciado.
    \item Visualizacion de precio, tiempo de produccion y ganancia.
    \item Compra de consumibles (pesticida y medicina).
    \item Control de monedas insuficientes.
    \item Integracion con inventario para la posterior colocacion.
\end{itemize}

\subsection{Colocacion, retiro y reutilizacion}
\begin{itemize}
    \item Colocacion manual en celda seleccionada por el usuario.
    \item Retiro de elementos del tablero hacia inventario.
    \item Recolocacion desde inventario en otra celda valida.
    \item Nombrado de animales al momento de colocar.
\end{itemize}

\subsection{Produccion y recoleccion}
\begin{itemize}
    \item Temporizador por elemento productivo.
    \item Opcion de recoleccion manual o automatica activable/desactivable.
    \item Conteo de tiempo restante visible en celdas.
    \item Suma automatica de ganancias al recolectar.
\end{itemize}

\subsection{Reglas de negocio del ciclo productivo}
\begin{itemize}
    \item Plantaciones: ciclo one-shot (comprar, sembrar, cosechar, volver a comprar).
    \item Animales y arboles: se mantienen en el terreno y reinician produccion tras recolectar.
    \item Recoleccion manual tardia: la produccion puede perderse por dano al exceder el limite de espera.
\end{itemize}

\subsection{Plagas y enfermedades}
\begin{itemize}
    \item Eventos aleatorios durante el ciclo de produccion.
    \item Plantas/arboles requieren pesticida.
    \item Animales requieren medicina.
    \item Si no se cura a tiempo, el elemento muere y deja de producir.
\end{itemize}

\subsection{Panel de informacion, bitacora y estadisticas}
\begin{itemize}
    \item Panel superior con jugador, monedas, terreno actual, elementos, listos y tiempo jugado.
    \item Bitacora de eventos importantes (compra, cosecha, curacion, muerte, guardado, etc.).
    \item Estadisticas acumuladas por tipo de producto y ganancias totales.
\end{itemize}

\section{Estructuras de datos utilizadas}
\subsection{Matrices}
Cada terreno se modela con una matriz 10x10 donde cada posicion contiene \texttt{None} (celda vacia) o un identificador de elemento colocado.

\subsection{Diccionarios principales}
\begin{itemize}
    \item \textbf{Catalogo de productos}: definiciones de precio, tiempo y ganancia.
    \item \textbf{Items colocados}: acceso por id unico para consultar estado en tiempo real.
    \item \textbf{Inventario}: cantidad por clave de producto o consumible.
    \item \textbf{Estadisticas}: compras y ganancias por tipo.
\end{itemize}

\subsection{Listas}
\begin{itemize}
    \item Lista de terrenos de la partida.
    \item Lista de entradas de bitacora.
\end{itemize}

\subsection{Objetos y dataclasses}
Se utilizan clases y dataclasses para estado estructurado: \texttt{GameState}, \texttt{Terrain}, \texttt{PlacedItem}, \texttt{Statistics}, entre otras.

\section{Documentacion interna y externa}
\subsection{Documentacion interna (codigo fuente)}
\begin{itemize}
    \item Comentarios explicativos en bloques clave de logica (temporizadores, hilos, validaciones, persistencia).
    \item Nombres de clases y metodos orientados a comprension del flujo.
    \item Separacion por secciones: modelo, persistencia, motor de juego e interfaz.
\end{itemize}

\subsection{Documentacion externa (este documento)}
\begin{itemize}
    \item Descripcion general de arquitectura y funcionamiento.
    \item Manual de uso paso a paso.
    \item Explicacion del manejo de archivos y estructura de datos.
    \item Criterios de mantenimiento y mejoras futuras.
\end{itemize}

\section{Manual de usuario}
\subsection{Requisitos para ejecutar}
\begin{enumerate}
    \item Tener Python 3 instalado.
    \item Instalar dependencias de \texttt{requirements.txt}.
    \item Entorno virtual opcional (no obligatorio para ejecutar el proyecto).
    \item Ejecutar desde la raiz del proyecto con \texttt{python main.py}.
\end{enumerate}

\subsection{Flujo basico de uso}
\begin{enumerate}
    \item En el menu principal, elegir \textit{Iniciar Partida} o \textit{Continuar Partida}.
    \item Abrir \textit{Tienda} y comprar un producto.
    \item Colocar el producto en una celda vacia.
    \item Esperar el tiempo de produccion indicado.
    \item Recolectar manualmente o activar \textit{Auto recoleccion}.
\end{enumerate}

\subsection{Uso del inventario}
\begin{enumerate}
    \item Comprar en tienda (el item se agrega al inventario).
    \item Abrir la ventana \textit{Inventario}.
    \item Elegir el item y presionar \textit{Colocar item seleccionado}.
    \item Hacer clic en una celda libre para ubicarlo.
\end{enumerate}

\subsection{Curacion de plagas/enfermedades}
\begin{enumerate}
    \item Comprar pesticida o medicina en tienda.
    \item Seleccionar la celda afectada.
    \item Presionar \textit{Curar} antes de que termine el contador de la afeccion.
\end{enumerate}

\subsection{Compra y navegacion entre terrenos}
\begin{enumerate}
    \item Presionar \textit{Comprar terreno (20000)} cuando haya fondos.
    \item Navegar con \textit{Terreno -}, \textit{Terreno +} o \textit{Ir a terreno}.
\end{enumerate}

\subsection{Guardado y salida}
\begin{itemize}
    \item El juego guarda automaticamente cada 30 segundos.
    \item Tambien puede guardarse manualmente con el boton \textit{Guardar}.
    \item Al cerrar ventana de juego, se guarda estado y se detienen hilos.
\end{itemize}

\section{Manejo de archivos}
\subsection{Archivo principal de persistencia}
El archivo \texttt{data/savegame.json} conserva el estado completo de la partida:
\begin{itemize}
    \item Jugador y monedas.
    \item Terrenos y ocupacion de celdas.
    \item Items colocados, tiempos de produccion y estados especiales.
    \item Inventario, bitacora y estadisticas.
    \item Configuracion de auto-recoleccion.
\end{itemize}

\subsection{Politica de tiempo y guardado}
\begin{itemize}
    \item El tiempo de produccion avanza solo mientras la app esta abierta.
    \item No se aplica progreso offline al reabrir el juego.
    \item Si se inicia partida nueva y se confirma, se elimina el archivo anterior.
\end{itemize}

\subsection{Recursos graficos}
\begin{itemize}
    \item Imagenes de productos y logo en \texttt{assets/images}.
    \item Fondo de menu configurable con \texttt{menu.jpg} (o alternos definidos en codigo).
\end{itemize}

\section{Cumplimiento de requisitos del enunciado}
\begin{longtable}{|p{0.36\textwidth}|p{0.56\textwidth}|}
\hline
\textbf{Requisito} & \textbf{Implementacion en el sistema} \\
\hline
POO (herencia y polimorfismo) & Jerarquia \texttt{FarmEntity} y subclases; fabrica polimorfica para logica de cura. \\
\hline
Interfaz Tkinter funcional & Menu, juego principal, tienda, inventario y estadisticas con interaccion completa. \\
\hline
Matrices para terrenos & Cada terreno se modela como matriz 10x10. \\
\hline
Manejo de archivos & Guardado y carga en JSON con estado integral de partida. \\
\hline
Hilos y temporizadores & Un hilo por item productivo para actualizar ciclos por segundo. \\
\hline
Inventario y retiro del tablero & Boton de retiro al inventario y recolocacion posterior desde ventana de inventario. \\
\hline
Plagas/enfermedades + cura & Eventos aleatorios con consumibles de cura y muerte por no atender a tiempo. \\
\hline
Navegacion entre terrenos & Botones de anterior/siguiente e ir directo por numero. \\
\hline
Estadisticas y bitacora & Acumulados de compras/ganancias y registro cronologico de eventos. \\
\hline
\end{longtable}

\section{Conclusiones}
El proyecto cumple de forma integral con los componentes solicitados: interfaz grafica completa, uso de matrices, persistencia JSON, POO con herencia/polimorfismo, hilos por produccion y flujo de juego funcional con economia, inventario, estadisticas, plagas/enfermedades y expansion de terrenos.

Desde la perspectiva de mantenimiento, la separacion por capas (modelo, motor, persistencia e interfaz) facilita extensiones futuras. Entre mejoras recomendadas se incluyen: niveles de dificultad, mas eventos aleatorios, mejores efectos visuales y exportacion de reportes de estadisticas.

\vspace{0.5cm}
\noindent\textbf{Autora:} Joselyn Hidalgo Torres

\end{document}
